\renewcommand{\tSoundText}{PSG del MSX para música y efectos.}
\renewcommand{\tVideoText}{\textbf{V9990} obligatorio para presentación, créditos, menú y juego.}
\renewcommand{\tMonitorText}{Acceso a las salidas de vídeo del MSX y del V9990: dos pantallas o un monitor con selección de entrada.}
\renewcommand{\tControlsText}{Cursores; Espacio o M para saltar. Joystick en cualquiera de los dos puertos, con cualquiera de sus botones.}
\renewcommand{\tSetMonitor}{Conecta ambas salidas de vídeo. Empieza por la del MSX para comprobar la detección y cambia al V9990 para el resto del juego.}
\renewcommand{\tSetOn}{Enciende el MSX. Su salida muestra \texttt{V9990 found!}. Si aparece \texttt{V9990 NOT found!}, apaga el equipo y comprueba el cartucho V9990.}
\renewcommand{\tBlackScreen}{\textbf{Si no ves la imagen}, revisa la entrada de vídeo (sección 3), los cartuchos y la RAM. El lanzador openMSX cambia de salida automáticamente.}
\renewcommand{\tCreditsTitle}{9. Créditos y derechos}
\renewcommand{\tSourceNote}{\href{https://gitlab.com/dihalt/madmix2-v9990}{Código fuente disponible en GitLab.} \textbf{\textexclamdown{}Mejora el juego si te apetece!}}
\newcommand{\tCreditsExtra}{\item Logo y sintonía inicial de DiHalt Studio: materiales de presentación de DiHalt Studio.\item Ilustración de apertura: arte aportado para esta edición. Logo del título, personajes y música del juego: materiales originales.}
\renewcommand{\tSetJoy}{Conecta un joystick al puerto 1 o 2 si quieres usarlo.}
\renewcommand{\tSetGame}{Inserta el cartucho del juego en la otra ranura. En MegaFlashROM SD SCC+, copia la ROM a la SD como \texttt{MADMIX2.ROM} y graba con \texttt{OPFXSD MADMIX2.ROM /A16}. \texttt{/U} es opcional: sin él, el juego usa un subslot; con él, usa el slot primario. \texttt{/A16} fuerza ASCII16 si no lo detecta automáticamente.}
